% from ega1-6-fr.tex
\oldpage[I]{150}
le rôle de «~paramètres~», ce qui donne à ces morphismes une signification
«~géométrique~». Les morphismes qui ne sont pas de type fini interviendront
surtout par la suite dans les questions de «~changement du préschéma de
base~», par exemple par localisation ou complétion.
\end{env}

\subsection{Détermination locale d'un morphisme.}
\label{subsection:I.6.5-fr}

\begin{proposition}[6.5.1]
\label{I.6.5.1-fr}
Soient $X$, $Y$ deux $S$-préschémas, $Y$ étant de type fini sur $S$~; soient
$x\in X$, $y\in Y$ au-dessus d'un même point $s\in S$.
\begin{enumerate}
  \item[(i)] Si deux $S$-morphismes $f=(\psi,\theta)$,
    $f'=(\psi',\theta')$ de $X$ dans $Y$ sont tels que
    $\psi(x)=\psi'(x)=y$, et que les $\mathscr{O}_s$-homomorphismes (locaux)
    $\theta_x^\sharp$ et ${\theta'}_x^\sharp$ de $\mathscr{O}_y$ dans
    $\mathscr{O}_x$ soient identiques, $f$ et $f'$ coïncident dans un
    voisinage ouvert de $x$.
  \item[(ii)] Supposons en outre $S$ localement noethérien. Pour tout
    $\mathscr{O}_s$-homomorphisme local
    $\varphi:\mathscr{O}_y\to\mathscr{O}_x$, il existe un voisinage ouvert
    $U$ de $x$ dans $X$ et un $S$-morphisme $f=(\psi,\theta)$ de $U$ dans $Y$
    tels que $\psi(x)=y$ et $\theta_x^\sharp=\varphi$.
\end{enumerate}
\end{proposition}

\begin{enumerate}
  \item[(i)] La question étant locale sur $S$, $X$ et $Y$, on peut supposer
    $S$, $X$, $Y$ affines d'anneaux respectifs $A$, $B$, $C$, $f$ et $f'$
    étant de la forme $({}^a\varphi,\widetilde{\varphi})$ et
    $({}^a\varphi',\widetilde{\varphi}')$ respectivement, où $\varphi$ et
    $\varphi'$ sont deux $A$-homomorphismes de $C$ dans $B$ tels que
    $\varphi^{-1}(\mathfrak{j}_x)={\varphi'}^{-1}(\mathfrak{j}_x)
    =\mathfrak{j}_y$, et les homomorphismes $\varphi_x$ et $\varphi'_x$ de
    $C_y$ dans $B_x$, déduits de $\varphi$ et $\varphi'$, sont identiques~;
    on peut en outre supposer que $C$ est une $A$-algèbre de type fini. Soient
    $c_i$ ($1\leq i\leq n$) des générateurs de la $A$-algèbre $C$, et posons
    $b_i=\varphi(c_i)$, $b'_i=\varphi'(c_i)$~; par hypothèse, on a
    $b_i/1=b'_i/1$ dans l'anneau de fractions $B_x$ ($1\leq i\leq n$). Cela
    signifie qu'il existe des éléments $s_i\in B-\mathfrak{j}_x$ tels que
    $s_i(b_i-b'_i)=0$ pour $1\leq i\leq n$, et on peut évidemment supposer
    tous les $s_i$ égaux à un même élément $g\in B-\mathfrak{j}_x$. On en
    conclut que l'on a $b_i/1=b'_i/1$ pour $1\leq i\leq n$ dans l'anneau de
    fractions $B_g$~; si $i_g$ est l'homomorphisme canonique $B\to B_g$, on a
    par suite $i_g\circ\varphi=i_g\circ\varphi'$~; donc les restrictions de
    $f$ et $f'$ à $D(g)$ sont identiques.
  \item[(ii)] On peut se ramener à la même situation que dans (i), et supposer
    en outre que l'anneau $A$ est noethérien. Soient $c_i$
    ($1\leq i\leq n$) des générateurs de la $A$-algèbre $C$, et soit
    $\alpha:A[X_1,\ldots,X_n]\to C$ l'homomorphisme de l'algèbre de polynômes
    $A[X_1,\ldots,X_n]$ sur $C$ transformant $X_i$ en $c_i$ pour
    $1\leq i\leq n$. Soit d'autre part $i_y$ l'homomorphisme canonique
    $C\to C_y$, et considérons l'homomorphisme composé
    \[
      \beta:A[X_1,\ldots,X_n]
      \xrightarrow{\alpha} C
      \xrightarrow{i_y} C_y
      \xrightarrow{\varphi} B_x.
    \]
    Désignons par $\mathfrak{a}$ le noyau de $\beta$~; comme $A$ est
    noethérien, il en est de même de $A[X_1,\ldots,X_n]$, et par suite
    $\mathfrak{a}$ admet un système fini de générateurs
    $Q_j(X_1,\ldots,X_n)$ ($1\leq j\leq m$). D'autre part, chacun des
    éléments $\varphi(i_y(c_i))$ peut s'écrire $b_i/s_i$, où $b_i\in B$ et
    $s_i\notin\mathfrak{j}_x$~; on peut en outre supposer tous les $s_i$ égaux
    à un même élément $g\in B-\mathfrak{j}_x$. Cela étant, on a par hypothèse
    $Q_j(b_1/g,\ldots,b_n/g)=0$ dans $B_x$~; posons
    \[
      Q_j(X_1/T,\ldots,X_n/T)
      =P_j(X_1,\ldots,X_n,T)/T^{k_j}
    \]
    où $P_j$ est homogène de degré $k_j$. Soit alors
    $d_j=P_j(b_1,\ldots,b_n,g)\in B$. Par hypothèse, on a $t_jd_j=0$ pour un
    $t_j\in B-\mathfrak{j}_x$ ($1\leq j\leq m$), et on peut évidemment
    supposer tous les $t_j$
